\textbf{UE :} S-INFO-015 Projet d'analyse et de conception logicielle\\
\textbf{Année académique :} 2026-2027\\
\textbf{Groupe :} G02 \textbf{Responsable de l'extension :} Quarta Hugo
(Matricule : 250711)

\hypertarget{gestion-du-budget-financier}{%
\subsubsection{Gestion du budget financier
:}\label{gestion-du-budget-financier}}

\hypertarget{feature-virements-bancaires-ruxe9currents}{%
\paragraph{1. Feature : Virements bancaires
récurrents}\label{feature-virements-bancaires-ruxe9currents}}

\hypertarget{us1---cruxe9ation-et-planification-dun-virement-ruxe9current}{%
\subparagraph{US1 - Création et planification d'un virement
récurrent}\label{us1---cruxe9ation-et-planification-dun-virement-ruxe9current}}

En tant que titulaire bénéficiant de l'autorisation de virement sur mon
portefeuille,\\
je veux programmer un virement bancaire récurrent en définissant sa
fréquence, son montant et ses comptes associés afin d'automatiser mes
transferts d'argent réguliers.

\textbf{Critères d'acceptation :}

\begin{itemize}
\tightlist
\item
  L'utilisateur sélectionne un compte source et un compte bénéficiaire
  (compte interne au portefeuille ou IBAN SEPA externe valide).
\item
  La fonctionnalité est accessible uniquement si la permission de
  virement a été accordée par l'institution financière.
\item
  L'utilisateur doit définir la période du virement : hebdomadaire,
  mensuelle ou trimestrielle.
\item
  L'utilisateur peut renseigner une date de début obligatoire et une
  date de fin optionnelle.
\item
  L'utilisateur peut associer une communication libre ou structurée et
  une catégorie budgétaire spécifique au virement.
\item
  Le virement programmé est visible dans la liste des virements
  récurrents avec le statut ``Actif''.
\item
  À chaque échéance, le montant est débité du compte source et crédité
  sur le compte bénéficiaire (et la transaction est enregistrée dans
  l'historique).
\item
  Si, à l'échéance, le solde est insuffisant, le virement récurrent
  échoue, aucune transaction n'est effectuée et une notification d'échec
  est envoyée au titulaire.
\end{itemize}

\hypertarget{us2---virements-ruxe9currents-conditionnels}{%
\subparagraph{US2 - Virements récurrents conditionnels
:}\label{us2---virements-ruxe9currents-conditionnels}}

En tant que titulaire bénéficiant de l'autorisation de virement sur mon
portefeuille,\\
je veux ajouter des conditions à mes virements récurrents afin
d'empêcher mon compte courant d'être à découvert ou d'épargner prudement
de l'argent.

\textbf{Critères d'acceptation :}

\begin{itemize}
\tightlist
\item
  Lors de la création ou modification d'un virement récurrent,
  l'utilisateur peut activer une des conditions suivantes :

  \begin{itemize}
  \tightlist
  \item
    Condition d'épargne prudente : Le virement du compte courant vers le
    compte d'épargne ne s'exécute que si le solde restant sur le compte
    courant après exécution reste supérieur (ou égale) à un seuil défini
    par l'utilisateur.
  \item
    Condition de compte à découvert : Un virement automatique est
    déclenché depuis un compte d'épargne vers le compte courant dès que
    le solde de celui-ci descend sous un seuil critique défini afin de
    le réapprovisionner.
  \end{itemize}
\item
  Si la condition n'est pas remplie à l'échéance, le virement est ignoré
  pour cette échéance sans être annulé pour les échéances futures et une
  notification d'information est envoyée à l'utilisateur.
\item
  L'utilisateur peut ajouter, modifier ou retirer une condition à
  n'importe quel instant sans devoir recréer le virement récurrent.
\end{itemize}

\hypertarget{us3---pause-et-ruxe9activation-des-virements-ruxe9currents}{%
\subparagraph{US3 - Pause et réactivation des virements récurrents
:}\label{us3---pause-et-ruxe9activation-des-virements-ruxe9currents}}

En tant que titulaire bénéficiant de l'autorisation de virement sur mon
portefeuille, je veux pouvoir suspendre, réactiver et supprimer un
virement récurrent programmé,\\
afin de contrôler l'ensemble de mes transferts automatiques actifs
facilement.

\textbf{Critères d'acceptation :}

\begin{itemize}
\tightlist
\item
  L'interface affiche la liste de tous les virements automatiques
  (actifs et suspendus).
\item
  L'utilisateur peut cliquer sur ``Mettre en pause'' : le statut du
  virement devient immédiatement ``En pause'' et aucune transaction
  n'est effectuée lors des échéances tant qu'il est dans cet état.
\item
  L'utilisateur peut réactiver un virement suspendu à tout moment : les
  exécutions reprennent alors à partir de la prochaine échéance.
\item
  L'utilisateur peut supprimer définitivement un virement récurrent (une
  demande de confirmation est affichée avant la suppression).
\item
  La suppression d'un virement récurrent n'efface pas l'historique des
  transactions passées de celui-ci.
\end{itemize}

\hypertarget{feature-classification-des-transactions}{%
\paragraph{2. Feature : Classification des
transactions}\label{feature-classification-des-transactions}}

\hypertarget{us4---personnalisation-et-gestion-des-catuxe9gories-budguxe9taires}{%
\subparagraph{US4 - Personnalisation et gestion des catégories
budgétaires}\label{us4---personnalisation-et-gestion-des-catuxe9gories-budguxe9taires}}

En tant que titulaire connecté,\\
je veux gérer mes catégories budgétaires en modifiant les catégories
existantes ou en créant de nouvelles catégories\\
afin d'adapter le système de gestion de budget à mes habitudes de
consommation personnelles.

\textbf{Critères d'acceptation :}

\begin{itemize}
\tightlist
\item
  L'application fournit par défaut des catégories prédéfinies :
  \emph{Santé, Travail, Loisirs, Sport, Énergie, Transport, Études,
  Autres}.
\item
  Chaque catégorie possède un nom unique, une icône et un code couleur.
\item
  L'utilisateur peut ajouter une nouvelle catégorie personnalisée en
  définissant son nom, son icône et sa couleur.
\item
  L'utilisateur peut modifier le nom, la couleur ou l'icône d'une
  catégorie existante (les modifications se répercutent dynamiquement
  sur toutes les transactions et graphiques associés).
\item
  L'utilisateur peut supprimer une catégorie personnalisée (les
  transactions qui y étaient associées basculent dans la catégories
  ``Autres'' ou dans une catégorie choisie par l'utilisateur lors de la
  suppression).
\item
  La catégorie système ``Autres'' ne peut être ni renommée ni supprimée.
\end{itemize}

\hypertarget{us5---classification-manuelle-des-transactions}{%
\subparagraph{US5 - Classification manuelle des transactions
:}\label{us5---classification-manuelle-des-transactions}}

En tant que titulaire connecté,\\
je veux affecter ou modifier manuellement la catégorie associée à une
transaction financière passée ou récente afin de modifier ou préciser la
répartition de mes dépenses et revenus.

\textbf{Critères d'acceptation :}

\begin{itemize}
\tightlist
\item
  Dans l'historique des transactions d'un compte ou d'un portefeuille,
  chaque ligne affiche la catégorie actuellement associée avec son icône
  colorée.
\item
  L'utilisateur peut cliquer sur la catégorie d'une transaction afin de
  lui attribuer une autre catégorie.
\item
  La mise à jour modifie dynamiquement les visualisations budgétaires.
\item
  Par défaut, toute les transactions qui ne correspondent à aucune règle
  automatique est rattachée à la catégorie ``Autres''.
\end{itemize}

\hypertarget{us6---duxe9finition-de-ruxe8gles-de-classification-automatique}{%
\subparagraph{US6 - Définition de règles de classification
automatique}\label{us6---duxe9finition-de-ruxe8gles-de-classification-automatique}}

En tant que titulaire d'un portefeuille,\\
je veux configurer des règles d'assignation automatique de catégories
afin de catégoriser automatiquement mes dépenses et revenus.

\textbf{Critères d'acceptation :}

\begin{itemize}
\tightlist
\item
  L'utilisateur peut créer une règle de classification automatique liée
  à son portefeuille en précisant :

  \begin{itemize}
  \tightlist
  \item
    L'attribut de la transaction à analyser (ex : IBAN
    émetteur/destinataire, communication ou montant).
  \item
    L'opérateur de comparaison (\emph{égal à}, \emph{supérieur à},
    \emph{inférieur à}, \emph{contient},\ldots).
  \item
    La valeur cible (nombre à 11 chiffres, chaîne de caractères ou
    montant en fonction du type d'attribut comparé).
  \item
    La catégorie cible à assigner automatiquement.
  \end{itemize}
\item
  L'utilisateur peut organiser les règles de classification automatiques
  définies par ordre de priorité.
\item
  Lorsqu'une nouvelle transaction est effectuée, les règles sont
  évaluées dans l'ordre de priorité défini par l'utilisateur et dès
  qu'une règle est valide, la catégorie correspondante est immédiatement
  assignée à la transaction.
\item
  L'utilisateur peut appliquer une nouvelle règle créée à l'historique
  complet de ses transactions.
\end{itemize}

\hypertarget{us7---consultation-et-filtrage-de-lhistorique-des-transactions-uxe9tendus-aux-catuxe9gories-extension-de-lus18-de-base}{%
\subparagraph{US7 - Consultation et filtrage de l'historique des
transactions étendus aux catégories (Extension de l'US18 de
base)}\label{us7---consultation-et-filtrage-de-lhistorique-des-transactions-uxe9tendus-aux-catuxe9gories-extension-de-lus18-de-base}}

En tant que titulaire connecté,\\
je veux visualiser la répartition de mes dépenses et de mes revenus par
catégorie sous forme de graphiques, de listes et de tableaux,\\
afin d'identifier les catégories de dépenses et de revenus les plus
importantes de mon budget.

\textbf{Critères d'acceptation :}

\begin{itemize}
\tightlist
\item
  L'interface d'analyse comporte un diagramme circulaire et un graphique
  à barres illustrant la répartition en pourcentage et en montant par
  catégorie.
\item
  L'utilisateur peut filtrer la visualisation par période, par catégorie
  budgétaire et par comptes spécifiques au sein d'un ou plusieurs
  portefeuilles.
\item
  L'utilisateur peut exporter les données visualisées filtrées par
  catégorie aux formats CSV et JSON (extension de l'US19 de base).
\end{itemize}

\hypertarget{feature-statistiques-budguxe9taires-et-analyse-de-tendances}{%
\paragraph{3. Feature : Statistiques budgétaires et analyse de
tendances}\label{feature-statistiques-budguxe9taires-et-analyse-de-tendances}}

\hypertarget{us8---consultation-des-statistiques-budguxe9taires}{%
\subparagraph{US8 - Consultation des statistiques
budgétaires}\label{us8---consultation-des-statistiques-budguxe9taires}}

En tant que titulaire connecté,\\
je veux consulter des statistiques détaillées sur mes dépenses et
revenus afin d'évaluer la distribution et la variabilité de mes dépenses
sur une période choisie.

\textbf{Critères d'acceptation :}

\begin{itemize}
\tightlist
\item
  L'interface dédiée aux statistiques affiche pour la période
  sélectionnée :

  \begin{itemize}
  \tightlist
  \item
    Le montant minimum et maximum.
  \item
    La moyenne et l'écart-type.
  \item
    La médiane, les quartiles et l'écart interquartile.
  \end{itemize}
\item
  L'utilisateur peut choisir de calculer ces statistiques sur l'ensemble
  des transactions ou de sélectionner/exclure certaines catégories.
\item
  L'utilisateur peut choisir la période sur laquelle les statistiques
  sont calculées.
\item
  Si aucune transaction n'existe pour les filtres appliqués, un message
  informatif est affiché.
\end{itemize}

\hypertarget{us9---visualisation-des-tendances-temporelles-et-de-la-saisonnalituxe9}{%
\subparagraph{US9 - Visualisation des tendances temporelles et de la
saisonnalité}\label{us9---visualisation-des-tendances-temporelles-et-de-la-saisonnalituxe9}}

En tant que titulaire connecté,\\
je veux analyser l'évolution temporelle de mon budget via des courbes
comparatives et des indicateurs de saisonnalité,\\
afin de visualiser les variations de mes dépenses et revenus d'une
période à l'autre.

\textbf{Critères d'acceptation :}

\begin{itemize}
\tightlist
\item
  L'application fournit un graphique linéaire comparant les dépenses
  d'une période par rapport à la période précédente.
\item
  L'interface affiche le taux de variation (en pourcentage et en euro)
  entre les deux périodes comparées.
\item
  L'utilisateur peut sélectionner et modifier l'ensemble des catégories
  à prendre en compte lors de la visualisation.
\item
  L'application fournit une vue mettant en évidence les dépenses/revenus
  majoritaires en fonction de la saison.
\end{itemize}

\hypertarget{feature-duxe9tection-danomalies-et-alertes}{%
\paragraph{4. Feature : Détection d'anomalies et
alertes}\label{feature-duxe9tection-danomalies-et-alertes}}

\hypertarget{us10---duxe9tection-danomalies-de-duxe9penses-par-rapport-uxe0-lhistorique-et-au-foyer-moyen}{%
\subparagraph{US10 - Détection d'anomalies de dépenses par rapport à
l'historique et au foyer
moyen}\label{us10---duxe9tection-danomalies-de-duxe9penses-par-rapport-uxe0-lhistorique-et-au-foyer-moyen}}

En tant que titulaire connecté,\\
je veux être notifié automatiquement dès qu'une dépense inhabituelle ou
anormale est détectée sur mon compte afin d'éviter les dépenses
déraisonnables et d'améliorer la gestion de mon budget financier.

\textbf{Critères d'acceptation :}

\begin{itemize}
\tightlist
\item
  Le système analyse en continu les transactions effectuée et détecte :

  \begin{itemize}
  \tightlist
  \item
    Les anomalie par rapport à l'historique personnel : Le système
    détecte les dépenses d'une catégorie qui dépassent significativement
    la moyenne habituelle du titulaire.
  \item
    Les anomalie par rapport au foyer moyen : Le système compare les
    dépenses d'une catégorie du titualire par rapport à celles du foyer
    moyen.
  \end{itemize}
\item
  Dès qu'une anomalie est détectée, une alerte est envoyée au titulaire
  en précisant la transaction concernée ainsi que le montant et l'écart
  constaté.
\item
  L'utilisateur a la possibilité de marquer une alerte comme
  ``Légitime'' (pour ne plus être alerté sur cet événement) ou de régler
  les seuils de détection des anomalies pour une cette catégorie.
\end{itemize}

\hypertarget{feature-pruxe9diction-et-projection-budguxe9taire}{%
\paragraph{5. Feature : Prédiction et projection
budgétaire}\label{feature-pruxe9diction-et-projection-budguxe9taire}}

\hypertarget{us11---projection-du-solde-dun-portefeuille}{%
\subparagraph{US11 - Projection du solde d'un
portefeuille}\label{us11---projection-du-solde-dun-portefeuille}}

En tant que titulaire connecté,\\
je veux visualiser une prédiction de l'évolution du solde de mes comptes
pour un certain portefeuille à une date ultérieure\\
afin d'anticiper mon budget future et anticiper d'éventuelles
difficultés financières.

\textbf{Critères d'acceptation :}

\begin{itemize}
\tightlist
\item
  L'utilisateur peut sélectionner une date cible de prédiction future.
\item
  L'application prédis le solde future du solde sur base :

  \begin{itemize}
  \tightlist
  \item
    des virements récurrents actifs sur la période.
  \item
    de la projection des dépenses et revenus sur la période estimée à
    partir des statistiques historiques et des tendances.
  \end{itemize}
\item
  L'utilisateur peut sélectionner un ensemble de catégories à exclure ou
  à prendre en compte lors de la projection.
\end{itemize}
